\chapter{Definición del proyecto}
Retomando lo dicho en la propuesta de solución, se llevó a cabo una aplicación móvil que sirve como herramienta para el personal de seguridad y el personal académico en temporadas de ETS para disminuir la suplantación de identidad dentro de la ESCOM. 

La aplicación móvil cuenta con varios apartados dependiendo del usuario que la esté usando, ya sea un alumno, un personal académico o un personal de seguridad, sin embargo, la parte fundamental de esta aplicación es la del docente y la del alumno en la que pueden realizar el reconocimiento facial. El alumno puede hacer el reconocimiento facial a su persona para rectificar que no haya ningún problema con su pase de asistencia al momento de llegar a realizar su ETS. El personal académico cuenta con la parte de reconocimiento facial como una forma de pasarle la asistencia al ETS al alumno. 

Es importante mencionar que el pase de asistencia no se realiza únicamente con el reconocimiento facial, se puede realizar mediante otras formas, como con un escanéo a la credencial del alumno o más fácil, si el personal académico ya conoce al alumno de antemano, lo puede especificar así al momento de registrar la asistencia al alumno. 

Por parte del personal de seguridad, puede registrar el ingreso a la unidad académica (ESCOM), por medio del escaneo de la credencial del alumno. 

La aplicación móvil cuenta con más secciones, tales como la visualización de los datos de los ETS, el calendario escolar del IPN, una sección de mensajes para que se los alumnos se puedan comunicar con el personal académico y para que el personal académico se pueda comunicar tanto con los alumnos como con otro personal académico. 